\par\clearpage\SourceLeaf{306}\ArithmeticHead{221}{R}{}
\begin{center}\textbf{Inhaltsübersicht.}\end{center}
\begingroup
\fontsize{9}{12}\selectfont
\setlength{\tabcolsep}{3pt}
\newcommand{\TP}[1]{{\fontsize{8}{10}\selectfont\textls[350]{#1}}}
\newcommand{\CT}[4]{#1 & #2\ \TP{#3} & #4\\}
\noindent\begin{tabular}{@{}r p{.82\linewidth} r@{}}
\CT{I.}{Einleitendes und Übersicht über das benutzte Material}{..............}{S. 3}
\CT{II.}{Traditionen über die erste Jugendzeit bis zum Abgang auf die Universität Oktober 1795}{..}{9}
\CT{III.}{Die historische Jugendperiode: 1796 bis zur Doktorpromotion am 16. Juli 1799}{.....}{18}
\CT{}{a) Die ersten Studiensemester in Göttingen}{.................}{18}
\CT{}{b) Die lemniskatischen Funktionen in den Jahren 1796—99}{............}{30}
\CT{}{c) Die Theorie des arithmetisch-geometrischen Mittels (agM.) bis zum 30. Mai 1799}{...}{41}
\CT{}{d) Grundlagen der Analysis. Komplexe Grössen}{................}{44}
\CT{}{e) Die Inauguraldissertation}{.......................}{57}
\CT{IV.}{Die Jahre 1799 und 1800. AgM. Scheda Ac. Allgemeine elliptische Funktionen}{.....}{62}
\CT{}{a) Allgemeines. Das agM. nach dem 30. Mai 1799}{...............}{62}
\CT{}{b) Die allgemeine Theorie der elliptischen Funktionen}{.............}{72}
\CT{V.}{Die nachgelassene Abhandlung über das agM. Ausbau zur Theorie der Modulfunktion}{..}{84}
\CT{}{a) Einleitendes}{............................}{84}
\CT{}{b) Die nachgelassene Abhandlung zur Theorie des agM.}{.............}{88}
\CT{}{c) Theorie der Modulfunktion}{......................}{93}
\CT{}{d) Anwendungen. Zusammenfassung}{....................}{106}
\CT{}{e) Die letzten Tagebuchnotizen von 1800. Asymptotisches. Mittelpunktgleichung. Klassen-\newline anzahl}{..............................}{114}
\CT{VI.}{Die Jahre 1801 bis 1810}{.........................}{117}
\CT{}{a) Bis zur Berufung nach Göttingen 1807. Allgemeines. Zahlentheoretische Anwendungen\newline der elliptischen Funktionen}{......................}{117}
\CT{}{b) Das Jahr 1808. Beginn der Korrespondenz mit \textsc{Schumacher}}{..........}{125}
\CT{}{c) Die Jahre 1809, 1810. Der bilineare Algorithmus. Persönliches}{.........}{129}
\CT{VII.}{Die Jahre 1811, 1812. \textsc{Gauss}sche Reihe. Funktionentheorie}{............}{134}
\CT{}{a) Korrespondenz mit \textsc{Fr. W. Bessel} 1810—1812}{...............}{134}
\CT{}{b) Die Reihe \(F(\alpha,\beta,\gamma,x)\) bei \textsc{Euler} und \textsc{Gauss}. Die Begriffe Reihe und Konvergenz}{.}{136}
\CT{}{c) Analyse der Disquisitiones von 1812}{...................}{142}
\CT{}{d) Die nachgelassene Fortsetzung der Disquisitiones. Funktionentheoretisches}{.....}{148}
\CT{}{e) Die Tagebuchnotizen des Jahres 1812. \textsc{Gauss}sche Logarithmen}{.........}{154}
\CT{VIII.}{Die Jahre 1813 bis 1827}{.........................}{155}
\CT{}{a) Die 1813 bis 1816 veröffentlichten Abhandlungen}{..............}{155}
\CT{}{b) Das agM. in der Korrespondenz mit \textsc{Schumacher} (1816) und in der Determinatio attractionis}{}{166}
\CT{}{c) Konforme Abbildung. Geometrisches}{..................}{170}
\CT{}{d) Die Hundert Theoreme}{.......................}{181}
\end{tabular}
\par\clearpage\SourceLeaf{307}\ArithmeticHead{222}{L}{L. SCHLESINGER, ÜBER GAUSS' FUNKTIONENTHEORETISCHE ARBEITEN.}
\noindent\begin{tabular}{@{}r p{.82\linewidth} r@{}}
\CT{IX.}{Die Jahre 1827 bis 1843}{.........................}{S. 183}
\CT{}{a) \textsc{Abels} und \textsc{Jacobis} Entdeckungen. Logarithmisches Potential}{.........}{183}
\CT{}{b) \textsc{Harriots} Lehrsatz. Zwei algebraische Theoreme}{..............}{187}
\CT{}{c) Konforme Abbildung der Ellipse. Pentagramma mirificum}{...........}{192}
\CT{X.}{Zunehmendes Alter. Doktorjubiläum. Letzte Lebensjahre}{............}{199}
\CT{}{a) Jubiläumsschrift. Plan einer Abhandlung über Reihenkonvergenz}{........}{199}
\CT{}{b) Akademische Tätigkeit in den letzten Lebensjahren. \textsc{Riemann}}{.........}{206}
\CT{}{Sachregister}{...............................}{211}
\CT{}{Namenregister}{..............................}{214}
\CT{}{Berichtigungen zu der Abh. 2 (\textsc{Gauss}’ Werke Bd. X, 2)}{...............}{217}
\CT{}{Berichtigungen zu den Bänden III und X,1 der Werke}{...............}{219}
\CT{}{Inhaltsübersicht}{..............................}{221}
\end{tabular}
\par\bigskip\begin{center}\rule{.23\linewidth}{.3pt}\end{center}
\par\clearpage\endgroup
